\setcounter{deso}{2}
\begin{name}
	{\tenchude}
	{\tendethi}
	{\tentruong}
	{\thoigian}
\end{name}
\begin{bt}[1,5 điểm]~
\begin{enumerate}[1.]
    \item Tính giá trị biểu thức $A=\sqrt{(\sqrt{3}-1)^2}+4\sqrt{3}-\sqrt{75}$.
    \item Giải phương trình $1,5x+3=2(x-3)$.
    \item Giải bất phương trình $-2x+1 \le 5+4x$.
\end{enumerate}
\end{bt}

\begin{bt}[2,0 điểm]~
\begin{enumerate}[1.]
    \item Cho hàm số $y=ax^2$ với $a \ne 0$ có đồ thị đi qua điểm $M(-4;-4)$. Xác định hệ số $a$.
    \item Cho biểu thức $M=\dfrac{x-2\sqrt{x}}{x-4}+\dfrac{2\sqrt{x}}{2\sqrt{x}+4}$ và $N=\dfrac{\sqrt{x}-1}{\sqrt{x}}$
    \begin{enumerate}
        \item Rút gọn biểu thức $M$.
        \item Tìm tất cả giá trị nguyên của $x$ để $M.N$ nhận giá trị nguyên.
    \end{enumerate}
\end{enumerate}
\end{bt}

\begin{bt}[1,5 điểm]%[Đề tham khảo tuyển sinh vào 10 - Quận Gò Vấp - Đề số 1]%[9D2T3]%[TS10B7]
Để chuẩn bị cho hội trại 26 tháng 3, lớp 9A đi đặt may áo lớp. Giá mỗi áo nam là $120$ nghìn đồng,
mỗi áo nữ là $110$ nghìn đồng. Vì mua số lượng nhiều nên được giảm $10 \%$ trên tồng giá tiền do đó cả lớp trả số tiền tổng cộng là $4437$ nghìn đồng. Hỏi lớp $9$A có bao nhiêu bạn nam và bao nhiêu bạn nữ. Biết rằng sĩ số của lớp là $43$ học sinh.
\loigiai{
Gọi $x$ là số học sinh nam, $y$ là số học sinh nữ của lớp $9A$ .\\
Theo đề bài ta có hệ phương trình sau: $\left\{\begin{array}{l}
x+y=43\\
\left(120x+110y\right).90\%=4437
\end{array}\right.$ $\Leftrightarrow\left\{\begin{array}{l}
x=20\\
y=23
\end{array}\right.$ (thỏa điều kiện).\\
Vậy số học sinh nam là $20$ và số học sinh nữ là $23$ .
}
\end{bt}

\begin{bt}[1,0 điểm]%[Đề tham khảo tuyển sinh vào 10 - Quận 1 - Đề số 3]%[TS10B6]
Một xe bồn chở nước sạch cho một  khu chung cư có $200$ hộ dân. Mỗi đầu của bồn chứa nước là $2$ nửa hình cầu (có kích thước như hình vẽ). Bồn chứa đầy nước và lượng nước chia đều cho từng hộ dân. Tính xem mỗi hộ dân nhận được bao nhiêu lít nước sạch? (làm tròn đến chữ số thập phân thứ hai, lấy $\pi=3,14$)
\begin{center}
    \begin{tikzpicture}[>=stealth, line join=round, line cap=round, font=\footnotesize, scale=1,rotate=-90]
            \def\r{2}
            \def\h{5}
            \path (0,0) coordinate (O)
            (0:\r) coordinate (A)
            ($(A)+(90:\h)$) coordinate (A');
            \draw[dashed] (A) arc (0:180:{\r} and {\r/4})
            ($(A)+(90:\h)$) arc (0:180:{\r} and {\r/4});
            \draw (A) arc (0:-180:{\r} and {\r/4})
            ($(A)+(90:\h)$) arc (0:-180:{\r} and {\r/4})
            (A)--($(A)+(90:\h)$) ($(180:\r)$)coordinate (B)--++(90:\h)
            (A)arc(0:-180:\r)
            ($(A)+(90:\h)$)arc(0:180:\r)
            ;
            \draw[<->] ([xshift=-.5cm]B)--++(90:\h)node[midway,sloped,rotate=-90,fill=white]{$3{,}62$ m};
            \draw[<->]([yshift=2.3cm]A')--++(180:2*\r)node[midway,sloped,fill=white,rotate=-90]{$1{,}8$ m};
        \end{tikzpicture}
\end{center}
\loigiai{
Tổng thể tích của bồn nước $V=\dfrac{4}{3}\cdot 3{,}14\cdot 0{,}9^3 + 3{,}14\cdot 0{,}9^2\cdot 3{,}62=12{,}26$ m$^3 = 12\,260$ lít.\\
Khi đó mỗi hộ dân nhận được số lít là $12\,260:200=61{,}3$ lít.
}
\end{bt}

\begin{bt}[1,0 điểm]
Một tấm bìa hình tròn được chia thành 4 phần được đánh số 1, 2, 3, 4 và được gắn vào một trục quay có mũi tên cố định ở tâm. Bạn Minh quay tấm bìa, bạn Thanh gieo một con xúc xắc cân đối. Giả sữ mũi tên dừng lại ở hình quạt ghi số $a$ và số chấm xuất hiện trên xúc xắc là $b$. Tính xác suất của biến cố "Tích hai số $a$ và $b$ là một số không chia hết cho 2".
\end{bt}

\begin{bt}[3,0 điểm]%[Đề tham khảo tuyển sinh vào 10 - Quận Gò Vấp - Đề số 1]%[9H3-9]%[TS10K10]
Cho $(O;R)$ có dây cung $AB$ không đi qua tâm. Trên tia đối của tia $AB$ lấy điểm $M$ bất kỳ. Vẽ hai tiếp tuyến $MC$ và $MD$ với $(O)$ sao cho $D$ thuộc cung lớn $AB$ ($C$ và $D$ là hai tiếp điểm). Gọi $H$ là trung điểm của đoạn thẳng $AB$. Qua $A$ vẽ đường thằng song song với $MC$ và cắt $CD$ tại $K$, $BK$ cắt $MC$ tại $N$.
\begin{listEX}
\item Chứng minh: 5 điểm ${M}, {C}, {H}, {O}, {D}$ cùng thuộc một đường tròn và tứ giác ADHK nội tiếp.
\item Chứng minh: $N$ là trung điểm của $MC$.
\item Gọi $E$ là trung điểm của đoạn thẳng $BD$ và $F$ là chân đường vuông góc kẻ từ $E$ xuống cạnh $AD$. Chứng minh: Khi $M$ di chuyển trên tia đối của tia $AB$ thì điểm $F$ luôn nằm trên một đường tròn cố định có tâm là trung điểm của đoạn thẳng $OH$.
\end{listEX}
\loigiai{~
\begin{center}
\begin{tikzpicture}[scale=1,font=\footnotesize,line join=round,line cap=round, >=stealth]
\def\r{4}
\def\h{9}
\pgfmathsetmacro{\a}{sqrt(\h^2-\r^2)};
\pgfmathsetmacro{\go}{acos(\r/\h)}
\path [name path=c](0,0) circle (\r cm) ; %vẽ đường tròn và đặt tên
\path (0,0) coordinate (O) (40:4) coordinate (B);
\path (180-\go:\r) coordinate (C) -- ([turn]90:\a) coordinate (M)--(180+\go:\r) coordinate (D) ; %gán các điểm tạo ra tiếp tuyến bằng sử dụng phép quay
\draw [name path=mb] (M)--(B);
\path [name intersections={of=c and mb}] (intersection-1) coordinate (B) (intersection-2) coordinate (A)  ;
\path  ($(A)!0.5!(B)$) coordinate (H) ($(M)!0.5!(C)$) coordinate (N) ($(A)!2!(O)$) coordinate (S) ($(B)!0.5!(D)$) coordinate (E) ($(A)!(E)!(D)$)  coordinate (F)  ;
\path (intersection of B--N and C--D) coordinate (K)  (intersection of A--K and B--C) coordinate (Q) (intersection of E--F and B--S) coordinate (L); %gán các giao điểm
\draw (O)--(M)--(C)--(O) --(D)--(M)  (C)--(D)  (B)--(C) (N)--(B) (H)--(K) (A)--(Q) (A)--(O)--(S)--(B)--(D) --(A) (F)--(L);
\draw (0,0) circle (\r cm) ;
\foreach \x/\g in  {O/30,A/120,M/180,B/65,C/60,D/-90,H/-90,N/150,K/120,Q/90,S/-20,E/-90,F/190,L/-120}
{\fill[black] (\x)  circle(1.5pt) ($(\x)+(\g:3mm)$) node{$\x$};}; %đặt tên các điểm
\end{tikzpicture}
\end{center}
\begin{listEX}
\item
•Tứ giác $MCOD$ nội tiếp đường tròn đường kính $OM$ .\\
Tứ giác $MHOD$ nội tiếp đường tròn đường kính $OM$ .\\
$\Rightarrow $ 5 điểm $M,C,H,O,D$ cùng thuộc một đường tròn.\\
•Vì $M,C,H,D$ cùng thuộc một đường tròn nên tứ giác $MCHD$ nội tiếp.\\
$\Rightarrow\widehat{HDC}=\widehat{HMC}$\\
Mà $\widehat{HMC}=\widehat{HAK}{\rm}\left(AK{\rm{//}}MC\right)$ $\Rightarrow\widehat{HDC}=\widehat{HAK}$ $\Rightarrow $ Tứ giác $ADHK$ nội tiếp.
\item Chứng minh: $N$ là trung điểm của $MC$.
Gọi $Q$ là giao điểm của$AK$ và $BC$ .\\
Ta có: $\widehat{KHA}=\widehat{KDA}$ (tứ giác $ADHK$ nội tiếp)\\
Mà: $\widehat{CBA}=\widehat{CDA}$ (hai góc cùng chắn)\\
$\Rightarrow\widehat{KHA}=\widehat{CBA}$ $\Rightarrow HK\parallel BC$\\
$\Rightarrow K$ là trung điểm $AQ$ .\\
Mặt khác:$\dfrac{AK}{MN}=\dfrac{BK}{BN}=\dfrac{QK}{CN}$ , $AK=QK$\\
$\Rightarrow MN=CN$ hay $N$ là trung điểm $CM$ .
\item Chứng minh: Khi $M$ di chuyển trên tia đối của tia $AB$ thì điểm $F$ luôn nằm trên một đường tròn cố định có tâm là trung điểm của đoạn thẳng $OH$. \\
Kẻ đường kính $AS$ của $(O)$ .\\
Gọi $L$ là giao điểm $EF$ và $BS$ , $I$ là giao điểm $AL$ và $OH$ .\\
Xét $\Delta BDS$ có: $\left\{\begin{array}{l}
EL\parallel DS\left(\perp AD\right)\\
EB=ED
\end{array}\right.$ $\Rightarrow L$ là trung điểm của $BS$ .\\
Xét $\Delta ABS$ có: $\dfrac{HI}{BH}=\dfrac{OI}{KS}\left(=\dfrac{AI}{AL}\right)$ $\Rightarrow OI=IH$ $\Rightarrow I$ là trung điểm của $OH$ \\
Ta có: $ABLF$ nội tiếp đường tròn đường kính $AL$ (do $\widehat{ABL}=\widehat{AFL}=90^\circ $)\\
Vì $I$ là trung điểm $AL$ nên tứ giác $ABLF$ nội tiếp đường tròn tâm $I$\\
Vậy khi $M$ di chuyển trên tia đối của tia $AB$ thì điểm $F$ luôn nằm trên một đường tròn cố định có tâm là trung điểm của đoạn thẳng $OH$ .\\
\end{listEX}
}
\end{bt}